\documentclass[12pt]{report}
\usepackage[a4paper,margin=2.5cm]{geometry}
\usepackage[utf8]{inputenc}
\usepackage[spanish]{babel}
\usepackage{graphicx}
\usepackage{setspace}
\usepackage{titlesec}
\usepackage{ragged2e}
\usepackage{amsmath} 
\usepackage{amssymb}
\usepackage{hyperref}
\usepackage{url}
\usepackage{listings}
\usepackage{xcolor}
\usepackage{mdframed}
\usepackage{caption}

% --- Estética Premium para el Trabajo Final ---
\definecolor{darkblue}{rgb}{0.0, 0.0, 0.55}
\definecolor{codegreen}{rgb}{0,0.6,0}
\definecolor{codegray}{rgb}{0.5,0.5,0.5}
\definecolor{codepurple}{rgb}{0.58,0,0.82}
\definecolor{backcolour}{rgb}{0.97,0.97,0.95}

\titleformat{\chapter}
  {\normalfont\Huge\bfseries\color{darkblue}}
  {\thechapter.}
  {1em}
  {}

\lstdefinelanguage{Haskell}{
  keywords={data, type, where, deriving, import, qualified, newtype, case, of, do, get, put, return, lift, throw, module, let, in, if, then, else, instance, where, when, forall, not, null, head, span, isDigit, isAlpha, isAlphaNum, tail, dropWhile, lookup, Set, Map},
  keywordstyle=\color{blue}\bfseries,
  commentstyle=\color{gray},
  stringstyle=\color{red!70!black},
  basicstyle=\ttfamily\small,
  morecomment=[l]--,
  morestring=[b]"
}

\lstset{
  language=Haskell,
  backgroundcolor=\color{backcolour},
  commentstyle=\color{codegreen},
  keywordstyle=\color{magenta},
  numberstyle=\tiny\color{codegray},
  stringstyle=\color{codepurple},
  basicstyle=\ttfamily\footnotesize,
  breakatwhitespace=false,         
  breaklines=true,                 
  captionpos=b,                    
  keepspaces=true,                 
  numbers=left,                    
  numbersep=5pt,                  
  showspaces=false,                
  showstringspaces=false,
  showtabs=false,                  
  tabsize=2
}

\newcommand{\bnf}{::=}

\begin{document}

% --- PORTADA ---
\begin{titlepage}
    \centering
    {\Large \textsc{Universidad Nacional de Rosario}\par}
    \vspace{0.8cm}
    {\large \textsc{Facultad de Ciencias Exactas, Ingeniería y Agrimensura}\par}
    \vspace{1.8cm}
    {\Large Análisis de Lenguajes de Programación\par}
    \vspace{0.3cm}
    {\large Licenciatura en Ciencias de la Computación\par}
    \vspace{1.5cm}
    {\large \textsc{INFORME TÉCNICO DE TRABAJO PRÁCTICO FINAL}\par}
    \vspace{0.4cm}
    {\Huge \textbf{Intérprete y Motor de Optimización para Álgebra Relacional}\par}
    \vspace{2.5cm}
    {\large \textsc{Alumno: Steeman, Juan Ignacio \quad S-5613/8}\par}
    \vfill
    {\large \textbf{Cátedra}\par}
    {\large Mauro Jaskelioff - Cecilia Casolo\par}
    \vspace{1.5cm}
    {\large Rosario, Argentina \\ Marzo, 2026\par}
\end{titlepage}

\tableofcontents
\newpage

% ================================================================
\chapter{Introducción}

\section{Contexto y Objetivos}
El objetivo de este trabajo fue diseñar e implementar un intérprete interactivo y de procesamiento por lotes para el \textbf{Álgebra Relacional}. El Álgebra Relacional (AR) constituye el fundamento teórico de los sistemas de gestión de bases de datos relacionales, proporcionando un marco formal para la manipulación de datos mediante la teoría de conjuntos.

La arquitectura de procesamiento del programa sigue las fases clásicas de los compiladores e intérpretes interactivos (REPL):\\
$Input \rightarrow Lexer \rightarrow Parser (AST) \rightarrow Opt (Optimizer) \rightarrow Eval \rightarrow \text{Output (PrettyPrinter)}$.

Nuestra herramienta permite no solo la evaluación directa de consultas, sino también la definición de esquemas, la inserción dinámica de datos y la optimización simbólica de las expresiones para mejorar el rendimiento de cómputo sobre relaciones de gran escala.

\section{Operaciones Soportadas}
El intérprete soporta las siguientes operaciones elementales:
\begin{itemize}
    \item \textbf{Selección ($\sigma$):} Filtrado de tuplas basado en predicados booleanos.
    \item \textbf{Proyección ($\pi$):} Reducción de la dimensionalidad de la relación.
    \item \textbf{Unión ($\cup$):} Combinación de relaciones con esquemas compatibles.
    \item \textbf{Diferencia ($-$):} Eliminación de elementos presentes en el sustraendo.
    \item \textbf{Producto Cartesiano ($\times$):} Combinación exhaustiva de tuplas.
    \item \textbf{Renombre ($\rho$):} Modificación de etiquetas de atributos para evitar colisiones.
\end{itemize}

Asimismo, se implementaron de forma íntegra operaciones derivadas fundamentales:
\begin{itemize}
    \item \textbf{Intersección ($\cap$):} Elementos comunes mediante $\text{Rel}_1 - (\text{Rel}_1 - \text{Rel}_2)$.
    \item \textbf{Natural Join ($\bowtie$):} Enganche automático sobre atributos compartidos.
    \item \textbf{División ($\div$):} Operador complejo utilizado para cuantificación universal.
\end{itemize}

% ================================================================
\chapter{El Lenguaje y su Sintaxis}

\section{Modelado Teórico en Haskell}
Para modelar las relaciones de forma matemáticamente fiel, utilizamos la estructura definida en \texttt{Commons.hs}:

\begin{lstlisting}
data Relacion = R {
    nombre    :: String,
    atributos :: [(Atributo, Type)],
    tuplas    :: Set.Set Tupla
}
\end{lstlisting}

La elección de \texttt{Data.Set} para almacenar las tuplas no es trivial: nos permite delegar en la librería la gestión de la unicidad y el orden de los elementos, garantizando que una relación siempre sea un **conjunto** de datos y no un multiconjunto, tal como exige el álgebra relacional pura.

Internamente, una tupla se modela como un mapeo:
\begin{lstlisting}
type Tupla = Map.Map Atributo Valor
\end{lstlisting}
Esto asocia cada atributo (nombre de columna) con su valor, permitiendo un acceso semántico robusto e independiente del orden físico en la lista de atributos para operaciones como el cálculo de predicados.

\section{Sintaxis Abstracta (BNF)}
A continuación, formalizamos la gramática del lenguaje procesada por \textit{Happy}:

\begin{align*}
Atributo &\bnf String \\
NombreRel &\bnf String \\
Expr &\bnf ERelacion\ NombreRel \\
     &\mid ESeleccion\ Cond\ Expr \\
     &\mid EProyeccion\ [Atributo]\ Expr \\
     &\mid EUnion\ Expr\ Expr \\
     &\mid EDiff\ Expr\ Expr \\
     &\mid EProd\ Expr\ Expr \\
     &\mid ERenombre\ Atributo\ Atributo\ Expr \\
     &\mid EInterseccion\ Expr\ Expr \\
     &\mid ENaturalJoin\ Expr\ Expr \\
     &\mid EDiv\ Expr\ Expr 
\end{align*}

\vspace{0.3cm}
\textbf{Condiciones y Predicados ($\sigma$):}
\begin{align*}
Cond &\bnf PTrue \mid PFalse \\
     &\mid PEq\ Atributo\ Valor \mid PNeq\ Atributo\ Valor \\
     &\mid PLt\ Atributo\ Valor \mid PGt\ Atributo\ Valor \\
     &\mid PAttrEq\ Atributo\ Atributo \\
     &\mid PAnd\ Cond\ Cond \mid POr\ Cond\ Cond \\
     &\mid PNot\ Cond
\end{align*}

% ================================================================
\chapter{Semántica y Evaluación}

\section{Implementación de Operaciones Críticas}
El módulo \texttt{Operations.hs} encapsula la lógica matemática. Destacan las siguientes implementaciones:

\subsection{Join Natural ($\bowtie$)}
A diferencia del producto cartesiano, el Join Natural busca atributos comunes. Nuestra implementación utiliza la función \texttt{compatibles} para filtrar el producto de todas las tuplas:
\begin{lstlisting}
compatibles tupla0 tupla1 =
    all (\attr -> Map.lookup attr tupla0 == Map.lookup attr tupla1) comunes
\end{lstlisting}
Si las tuplas coinciden en todos sus atributos compartidos, se realiza la unión de los mapeos (\texttt{Map.union}).

\subsection{División Relacional ($\div$)}
La división es la operación más compleja. Se implementó siguiendo la definición formal que utiliza proyecciones y diferencias:
\begin{lstlisting}
division r s = do
    proy1 <- proyeccion atrib r
    let prod = productoCartesiano proy1 s
    diff  <- diferencia prod r
    proy2 <- proyeccion atrib diff
    diferencia proy1 proy2
\end{lstlisting}
Este enfoque garantiza que el resultado contenga solo aquellas tuplas de la proyección del dividendo que están relacionadas con **todas** las tuplas del divisor.

% ================================================================
\chapter{Infraestructura Monádica}

\section{El Transformador \texttt{StateError}}
Para gestionar el estado mutable (catálogo de relaciones) y el manejo de errores (excepciones semánticas), construimos un stack monádico en \texttt{Monads.hs}:

\begin{lstlisting}
type StateError a = StateT State (Either Error) a
\end{lstlisting}

Este diseño proporciona ventajas críticas:
\begin{enumerate}
    \item \textbf{Persistencia del Estado:} Incluso si una consulta falla (ej. error de tipos), el estado anterior del intérprete se conserva.
    \item \textbf{Errores Semánticos:} Permite capturar fallos como \texttt{RelacionNoExiste} o \texttt{AtributosNoCompatibles} sin interrumpir la ejecución del programa principal ni recurrir a \texttt{error} de Haskell.
    \item \textbf{Pureza:} La lógica de evaluación sigue siendo pura y fácil de testear, ya que los efectos están controlados por el sistema de tipos.
\end{enumerate}

% ================================================================
\chapter{Motor de Optimización Heurística}

\section{Estrategia de Push-down Selection}
La consulta $\sigma_{C}(R_1 \times R_2)$ es computacionalmente costosa si realizamos primero el producto cartesiano. Nuestro optimizador (\texttt{Optimizador.hs}) analiza el árbol y aplica la regla de "empujar" las selecciones:

Si los atributos requeridos por el predicado $C$ pertenecen únicamente a la relación $R_1$, la expresión se transforma en $(\sigma_{C} R_1) \times R_2$. Esto reduce el tamaño de las relaciones intermedias de forma exponencial, minimizando el uso de memoria.

\section{Otras Simplificaciones}
\begin{itemize}
    \item \textbf{Combinación de selecciones:} $\sigma_{C1}(\sigma_{C2}(R)) \rightarrow \sigma_{C1 \land C2}(R)$.
    \item \textbf{Eliminación de proyecciones redundantes:} $\pi_{A1}(\pi_{A2}(R)) \rightarrow \pi_{A1}(R)$ (siempre que $A1 \subseteq A2$).
\end{itemize}

% ================================================================
\chapter{Manual de Uso e Instalación}

\section{Compilación con Stack}
El proyecto utiliza \textit{Stack} para la gestión de dependencias y builds.
\begin{lstlisting}[language=bash]
# Compilar el proyecto
stack build

# Iniciar el interprete REPL
stack exec TP-AR-exe
\end{lstlisting}

\section{Guía de Comandos del Intérprete}
El REPL soporta comandos especiales iniciados por dos puntos:
\begin{itemize}
    \item \texttt{:createRel [Nombre] [Atrib1:Tipo, ...]} - Crea una nueva relación.
    \item \texttt{:insertRel [Nombre] (Val1, Val2), ...} - Inserta tuplas.
    \item \texttt{:load [Archivo]} - Carga un script relacional (\texttt{.ar}).
    \item \texttt{:browse} - Muestra las relaciones actualmente en memoria.
    \item \texttt{:help} - Muestra la ayuda interactiva.
\end{itemize}

\section{Ejemplo Completo de Consulta}
\begin{lstlisting}
// Definicion
:createRel Empleados id:int, nombre:string, dep:int;
:createRel Deps id:int, nombre:string;

// Consulta: Nombres de empleados del departamento "Ventas"
Ventas = seleccion[nombre = 'Ventas'](Deps);
EmpVentas = naturaljoin(Empleados, Ventas);
Resultado = proyeccion[nombre](EmpVentas);
\end{lstlisting}

% ================================================================
\chapter{Conclusiones}
La implementación en Haskell permitió abstraer la complejidad matemática del Álgebra Relacional mediante un sistema de tipos sólido. La combinación de \texttt{Happy} para el parseo y una infraestructura monádica robusta resultó en un intérprete estable, extensible y capaz de manejar optimizaciones de consultas de forma elegante.

% ================================================================
\begin{thebibliography}{9}
\bibitem{AR} Wikipedia. \textit{Relational Algebra}.
\bibitem{Codd} Codd, E. F. (1970). \textit{A Relational Model of Data for Large Shared Data Banks}.
\bibitem{Happy} Hackage. \textit{Happy Documentation}.
\end{thebibliography}

\end{document}
